\section{Introduction}

Retrieval-augmented generation (RAG) papers often report an improvement in a
retrieval score and then infer an improvement in answer faithfulness. The
reasoning is understandable. If the retriever returns passages that better
match the query, then the generator should have better evidence. Dense
retrieval, reranking, compression, and post-retrieval refinement all inherit
some version of this local premise.

The premise is incomplete. A generator receives a context set, not independent
passage scores. Individually relevant passages can be redundant, mutually
inconsistent, or missing the bridge needed to support an answer. Conversely, a
lower-scoring passage may contain the one fact that makes the answer grounded.
Faithfulness also depends on the evaluator: an NLI scorer, another NLI scorer,
and an LLM judge can assign materially different support values to the same
answer. A deployment claim then adds another layer: a method that improves one
dataset row may be too slow, too threshold-sensitive, or worse on another task.

The central question is therefore:

\begin{quote}
\emph{Which RAG faithfulness claims remain identifiable once local retrieval
similarity, context-set structure, evaluator choice, transfer, and cost are
separated?}
\end{quote}

We answer with a controlled falsification suite built around a tempting but
ultimately rejected mechanism. The starting variable is the Context Coherence
Score (CCS), a retrieval-time statistic that measures mean pairwise similarity
among retrieved passages. CCS is a plausible set-level complement to
query-passage relevance: it asks whether the retrieved passages look mutually
compatible, rather than whether each passage looks relevant in isolation. A
natural hypothesis is that this coherence variable explains why some retrieval
refinements appear to reduce faithfulness.

The controlled test rejects the strong version of that story. We construct
HIGH- and LOW-CCS context sets that are matched on mean query-passage
similarity. The manipulation creates a large coherence gap; the faithfulness
effect is null and slightly negative. This null is not a cosmetic limitation.
It is the central methodological result: an intuitive RAG mechanism can look
plausible under aggregate retrieval and faithfulness tables, then disappear
when the relevant controls are imposed.

The rest of the paper shows that this failure is part of a wider
identifiability problem. Local retrieval similarity does not identify answer
support. CCS alone does not identify answer support. A single automatic metric
does not identify intervention effect size. A threshold tuned on one dataset
does not cleanly identify a transferable policy. A one-row baseline comparison
does not identify deployment superiority once latency and indexing cost are
included.

\paragraph{Contributions.}
\begin{itemize}[leftmargin=1.6em,itemsep=2pt,topsep=2pt]
  \item \textbf{A controlled falsification of causal coherence.} We isolate
  context coherence from mean query similarity and show that higher CCS does
  not improve faithfulness in a preregistered SQuAD/Mistral intervention.
  \item \textbf{A scale correction of the refinement paradox.} The original
  pilot-sized DeBERTa effect shrinks to 1.1 points at $n=500$ across five
  seeds, with only one seed reaching significance for the baseline--HCPC-v1
  contrast.
  \item \textbf{A multi-metric non-identifiability result.} The same
  generations yield materially different conclusions under DeBERTa, a second
  NLI model, and a local RAGAS-style judge.
  \item \textbf{A threshold-transfer and baseline/cost audit.} We show that
  threshold recovery is dataset-dependent and that HCPC-v2, CRAG, RAPTOR-2L,
  and Self-RAG do not form a single deployment ranking.
  \item \textbf{A released controlled protocol and artifact.} We package
  matched controls, per-query generations, metric scores, threshold-transfer
  matrices, noise controls, and Pareto accounting as \method.
\end{itemize}

The paper's strongest claim is deliberately narrow: standard RAG reporting is
insufficient without controls that separate retrieval scores, context-set
structure, evaluator choice, transfer, and cost.
